\documentclass[prl,aps,preprint,floatfix]{revtex4-2}
\usepackage{graphicx}
\usepackage{amsmath,amssymb}
\usepackage{dsfont}
%\usepackage{hyperref}
\usepackage{xcolor}
% Custom definitions start
\newcommand{\Hq}{\mathcal{H}}
\newcommand{\heff}{H_{\scriptscriptstyle \rm eff}}
\newcommand{\hs}{\mathcal{H}_{\scriptscriptstyle \rm S}}
\newcommand{\hdr}{\mathcal{H}_{\scriptscriptstyle \rm DR}}
\newcommand{\Lhl}{\mathcal{H}_{\scriptscriptstyle \rm L}}
\newcommand{\hl}{\mathcal{H}_{\scriptscriptstyle \rm L}}
\newcommand{\ul}{U_{\scriptscriptstyle \rm L}}
\newcommand{\hsl}{\mathcal{H}_{\scriptscriptstyle \rm SL}}
\newcommand{\hint}{\mathcal{H}_{\scriptscriptstyle \rm int}}
\newcommand{\Lhsl}{\mathcal{H}_{\scriptscriptstyle \rm SL}}
\newcommand{\hdd}{\mathcal{H}_{\scriptscriptstyle \rm dd}}
\newcommand{\Lhloc}{\mathcal{H}_{\scriptscriptstyle \rm Loc}}
\newcommand{\hloc}{H_{\scriptscriptstyle \rm loc}}
\newcommand{\nn}{\nonumber}
\newcommand{\tq}{\mathcal{T}}
\newcommand{\fo}{\mathcal{F}}
\newcommand{\wo}{\omega_{\circ}}
\newcommand{\Ur}{U^{\dagger}_{L}}
\newcommand{\tr}{\text{Tr}}
\newcommand{\trl}{\text{Tr}_{ \scriptscriptstyle \rm L}}
\newcommand{\rh}{\rho_{\scriptscriptstyle  S}}
\newcommand{\rl}{\rho_{\scriptscriptstyle \rm L}^{\scriptscriptstyle \rm eq}}
%\newcommand{\sc}{ \scriptscriptstyle \rm sec}
\newcommand{\so}{\mathcal{S}}
\newcommand{\etal}{\textit{et al.\ }}
\newcommand{\sts}{\vert_{t\rightarrow\infty}}
\newcommand{\stoo}{\vert_{t\rightarrow 0}}
\newcommand{\ic}{\vert_{t=0}}
\newcommand{\ket}[1]{\vert #1 \rangle}
\newcommand{\bra}[1]{\langle #1 \vert}
\newcommand{\comment}[1]{ }
\begin{document}
\title{Thermal fluctuations and stability of discrete time-crystalline phases}
\author{Saptarshi Saha}
\email{ss17rs021@iiserkol.ac.in}
\author{Rangeet Bhattacharyya}
\email{rangeet@iiserkol.ac.in}
\affiliation{Department of Physical Sciences, Indian Institute of Science Education and Research Kolkata,
Mohanpur -- 741246, WB, India}

\begin{abstract}

%Discrete time crystal (DTC) is one of the recently-reported dynamical phases created by
%periodic drives on a quantum many body system (QMB). DTCs have been predicted and later
%experimentally demonstrated to show robust subharmonic responses. In this work, we treat
%the QMB as an open system and use a fluctuation-regulated quantum master equation (FRQME)
%to analyze its response to periodic drive. Motivated by the recent experiments on DTC, we
%choose a dipolar system subjected to external drive. FRQME captures the dissipative
%effects of the drive and dipolar coupling on the dynamics regularized by the thermal
%fluctuations. In this letter, we show that such terms lend stability to the dynamics and
%is directly responsible for the robustness.  Specifically, we find longer fluctuation
%correlation time enhances the stability of DTC. Our results are in good agreement with the
%experiments. We also show how temperature affects DTC performances.


%DTC in QMB systems are currently under intense investigation. An important milestone in
%the development is experimental realizations of DTC using open systems. DTCs show
%subharmonic responses and are robust.  Here, we present an alternate theoretical
%description using a FRQME. As our system, we choose a dipolar system subjected to external
%drives same as was in recent experiments. FRQME captures the dissipative effects of the
%drive and dipolar coupling on the dynamics regularized by the thermal fluctuations. In
%this letter, we show that such terms lend stability to the dynamics and is directly
%responsible for the robustness.  Specifically, we find longer fluctuation correlation time
%enhances the stability of DTC. Our results are in good agreement with the experiments. We
%also show how temperature affects DTC performances.

Breaking of the discrete time-translation symmetry leads to the emergence of the discrete
time-crystalline (DTC) phase in quantum many-body systems. In this phase, system
observables exhibit a robust sub-harmonic response. DTC has been experimentally realized
in driven dipolar system which are usually analyzed using a Floquet formalism
applicable for closed systems.
Here, we extend the analysis to a realistic open quantum
system, a dipolar coupled two-spin dissipative system subjected to a periodic drive. 
%The excitation protocol consists of two non-commuting Hamiltonian applied in successive time steps. 
To this end, we use a
fluctuation-regulated quantum master equation (FRQME) for our analysis. The dissipators of this
master equation are regularized by thermal fluctuations and play a central role in
stabilizing the DTC phase against perturbations. Our results are in excellent agreement
with the published experiments. Moreover, we show the temperature-dependence of the DTC
phase in open quantum systems.

%Breaking of the discrete time-translation symmetry leads to the emergence of the discrete
%time-crystalline (DTC) phase in quantum many-body systems. The system observables exhibit
%a robust sub-harmonic response in this phase.  Motivated by the recent experiments, we
%provide an alternative theoretical prescription to describe the origin of such
%sub-harmonic response in a dipolar coupled two-spin dissipative system subjected to a
%periodic drive that consists of two non-commuting Hamiltonian applied in successive time
%steps. To the end, we use a fluctuation-regulated quantum master equation for our
%analysis. We show that the second-order dissipative process plays a crucial role to
%stabilize the DTC phase under perturbations. We also obtain the phase diagram of DTC which
%matches with the existing experimental results. In addition to that, our result displays
%the dependency of the DTC phase on thermal fluctuation, which indicates that the stability
%of the DTC phase depends on the temperature of the open quantum system. 





\end{abstract}

\maketitle

Investigating the out-of-equilibrium phases of a quantum system is one of the emerging
research areas in physics \cite{Polkovnikov2011,eisert_quantum_2015}. There exist several
exotic phases of a many-body quantum system that are usually forbidden in thermal
equilibrium. Such phases perceive to be important in quantum computation and quantum
information processing \cite{abobeih_atomic-scale_2019, frey_realization_2022}.
Time-dependent periodic perturbation is proved as a useful tool for manipulating those
non-equilibrium phases \cite{Goldman2014,Abanin2015}. Some prime examples of such phases
are given as, many-body localization (MBL), prethermalization, and discrete
time-crystalline (DTC) \cite{abanin2017, ueda_2020,yao_discrete_2017}. Among them, the
existence of the DTC phase is relatively a newer topic.  Therefore, a large amount of
theoretical and experimental work is progressing in the DTC phase in the last few years.
\cite{sacha_time_2018, khemani2019}.

A space-crystal is formed when the continuous space-translation symmetry spontaneously breaks to a discrete
space-translation symmetry in the system \cite{ashcroft1976}. In 2012, Wilczek first proposed the concept of
breaking the continuous time translation symmetry \cite{wilczek_quantum_2012}. Although the realization of
such symmetry breaking is under big question, as in the presence of the time-periodic perturbation, one
cannot expect a time-independent many-body system to show time-periodicity in its equilibrium phase. Later,
a ``no-go theorem" is implemented to eliminate the possibility of having the continuous time-crystal (CTC)
\cite{bruno_impossibility_2013, watanabe_absence_2015}.

Motivated by the above proposal, Else \etal defined the DTC phase, which is characterized by a breaking of
the discrete time-translation symmetry in the presence of a time-periodic Hamiltonian
\cite{else_floquet_2016}. As such, the periodicity of any system observable is an integer multiple of the
period of the drive Hamiltonian. In most cases, a period-doubling response is observed, which is also
signified as a sub-harmonic response.

Later, several experiments using nuclear magnetic resonance (NMR), and trapped ions were designed to
implement the DTC phase \cite{zhang_observation_2017, choi_observation_2017, rovny_observation_2018,
pal_temporal_2018, kyprianidis_observation_2021, beatrez_critical_2023}.  One common feature for most of the
experiments mentioned above is the presence of an external drive that contains two or more non-commuting
Hamiltonians applied in successive time steps. As we are interested in the dynamics of a dissipative dipolar
system under periodic perturbation, we consider only the experimental demonstration by Choi \etal and
Beatrez \etal in this manuscript \cite{choi_observation_2017, beatrez_critical_2023}. 

Choi \etal showed the experimental observation of $2\tau$ periodic response of an ensemble of dipolar
coupled nitrogen-vacancy (NV) centers in a diamond. Here $\tau$ is the periodicity of the applied drive. The
drive consists of a spin-locking sequence along the `$x$' direction followed by a rotation along the `$y$'
axis. The emergence of the DTC phase is the consequence of a critically slow thermalization in the system.
Later Yao \etal provided a theoretical explanation of such experimental observations
\cite{ho_critical_2017}. Recently, using a similar kind of pulse sequence, Beatrez \etal observe a robust
sub-harmonic response in a dipolar coupled $^{13}$C nuclear spins in a diamond \cite{beatrez_critical_2023}.
The spin-locking pulse is used to lock the `$x$' magnetization in the prethermal
phase\cite{beatrez_floquet_2021}, which plays a significant role to stabilize the DTC phase under
perturbations. 

In general, for an isolated many-body system, the prethermalization in presence of periodic drive is
theoretically analyzed by  Floquet theory (i.e., Floquet prethermalization)
\cite{beatrez_floquet_2021,kuwahara_floquetmagnus_2016,bukov_prethermal_2015,rubio-abadal_floquet_2020,yin_prethermal_2021,holthaus_floquet_2016}.
The effect of the Floquet heating in each cycle can be captured by the second order contribution of the
periodic drive \cite{rubio-abadal_floquet_2020}. In the presence of local interaction and high-frequency
drive, such heating becomes exponentially slow, which results in a long-lived prethermal plateau in the
system \cite{Abanin2015}.  

 Although such systems are not perfectly isolated systems, they are very weakly coupled with the external
environment with a very-long relaxation time \cite{beatrez_floquet_2021,peng_floquet_2021}.  Recently Saha
\etal and Chakrabarti \etal provided a dynamical approach to describe the prethermalization in periodically
driven dissipative dipolar systems \cite{saha_cascaded_2023,chakrabarti_2022_creation}. We have used a
recently proposed fluctuation-regulated quantum master equation (FRQME) \cite{chakrabarti2018b}, which
successfully explains the interplay between the secular part of the on-resonant periodic drive and dipolar
interaction that leads to a prethermal phase. The dynamics is constrained by a set of quasi-conserved
quantities in this regime.  Subsequently, the non-secular interactions and system-bath interaction provide a
very slow thermalization process. Thus, the second-order terms of the drive and dipolar interaction
regulated by thermal fluctuations conveniently explain the effect of the Floquet heating in the system. 


The analytical form of the FRQME is given as \cite{chakrabarti2018b},
\begin{eqnarray} \label{frqme}
\frac{d\rh}{dt}&=& -i \trl\Big[\heff(t),\rh \otimes\rl\Big]^{\rm sec}\nn\\
&&-\int\limits^{\infty}_0 d\tau \trl\Big[\heff(t),\Big[\heff(t-\tau),\rh \otimes\rl\Big]\Big]^{\rm
sec}e^{-\frac{\vert\tau\vert}{\tau_c}},
\end{eqnarray} 
Here, $\rh$ is the system density matrix and $\rl$ is the equilibrium density matrix of the bath. It is
assumed that at the initial time, the system and bath are uncorrelated (Born approximation). `$\rm sec$'
denotes the secular approximation. $\heff$ is the interaction representation of $\hint + \hsl$. Where $\hsl$
represents the weak system-bath coupling and $\hint$ denotes the local interaction in the system (i.e.,
applied drive, dipolar interaction). $\tau_c$ is the fluctuation correlation time of the bath. The most
unique feature of the above Eq. (\ref{frqme}) is the presence of the exponential kernel in the second-order
terms, which is not present in the formal form of the Born-Markov QME. Such a memory kernel is necessary to
explain the dissipative effects coming from the local interactions (i.e., $\hint$) along with $\hsl$. FRQME
can be used as an essential tool for analyzing several applications in NMR, quantum optics, and quantum
information processing \cite{chanda2020,chatterjee_nonlinearity_2020,chanda2021,saha2022,saha_effects_2022}.


Here, we consider a dipolar coupled two-spin system. The system is also coupled with the bath. The dipolar
coupling Hamiltonian is denoted by $\hdd = \omega_d \left(\vec{I_1}.\vec{I_2} -
3(\vec{I_1}.\hat{r})(\vec{I_2}.\hat{r}) \right)$. $\omega_d$ is the dipolar coupling strength. and $I_{a}^i
= \sigma_{\alpha}^i/2, \,a \in \{x,y,z\}$, $\sigma_{a}$ is $a$th component of Pauli spin operators for
spin-1/2 particles. We assume that both the spins are identical. For a thermal bath consisting of harmonic
oscillators, $\hsl$ can be modeled by using the Jaynes-Cummins Hamiltonian \cite{breuer2002}. 
\begin{figure}[htb]
\includegraphics[width=0.64\linewidth]{Figure-1.pdf} 
\caption{The schematic diagram shows the experimental realization for demonstrating the DTC phase in the
dissipative dipolar system. The initial state is prepared by using a  $\pi/2$ pulse in the  `$y$' direction.
The two-drive protocol is given here. The spin-locking sequence is provided for $\tau_1$ time, secondly, a
$\pi + \delta $ rotation is given along `$y$' direction for a duration $\tau_2$. The whole sequence, $\tau =
\tau_1 + \tau_2$ is repeated for $n$ times before measuring the final spectrum of $S(M_x(t))$.}
\label{fig-1}
\end{figure}

Motivated by the two-pulse protocol by Choi and Beatrez \etal
\cite{choi_observation_2017,beatrez_critical_2023}, we use a similar kind of excitation scheme for our
theoretical calculations. Initially, the spins are oriented along the `$z$' direction in the presence of the
Zeeman field. A $\pi/2$ pulse along the `$y$' direction helps the system to orient along the  `$x$'
direction, so, $M_x \stoo =1$. In terms of wave function, we choose the initial state as, $\vert\psi\rangle
\stoo = \prod_{i=1}^2 \left(\vert\uparrow \rangle_i + \vert\downarrow \rangle_i\right)/ \sqrt{2}$.
$\vert\uparrow \rangle$ and $\vert\downarrow \rangle$ are the eigenstate of the Zeeman basis. Next, a
resonant periodic drive along the `$x$' direction is applied in the system for a time duration of $\tau_1$.
The corresponding Hamiltonian is given by, $\mathcal{H}_x(t) = \sum\limits_{i=1}^2 \omega_1 I_{x}^i
\cos\omega_\circ t$. We are interested in the non-equilibrium dynamics of the system in the presence of the
periodic drive, therefore we choose a system having weak system-bath coupling. The dynamics is dominated by
decay from other sources, i.e., dipolar dissipation and drive-induced dissipation in this case.

Using FRQME (Eq. (\ref{frqme})), the dynamical equation for this case is written as \cite{saha_cascaded_2023},
\begin{eqnarray}\label{dynamics-1}
\frac{d\rh}{dt}&=& \left( \mathcal{D}_{\scriptscriptstyle \rm sec} + \mathcal{D}_{\scriptscriptstyle \rm nsec} + \mathcal{D}_{\scriptscriptstyle \rm SL}\right) [\rh(t)]
%\frac{d\rh}{dt}&=& -i\left[H^{\scriptscriptstyle \rm sec}, \rh(t)\right] -  \tau_c\left[ H^{\scriptscriptstyle \rm sec},\left[H^{\scriptscriptstyle \rm sec},\rh(t) \right] \right],\qquad [0<t<\tau_1]
\end{eqnarray}
We argue that, the contribution from $\mathcal{D}_{\scriptscriptstyle \rm nsec}[\rh(t)]$ and $\mathcal{D}_{\scriptscriptstyle \rm SL}[\rh(t)]$ are relatively smaller than $\mathcal{D}_{\scriptscriptstyle \rm sec}[\rh(t)]$. The analytical form is given as,
\begin{eqnarray}
\mathcal{D}_{\scriptscriptstyle \rm sec}[\rh(t)] = -i\left[H_{\scriptscriptstyle \rm sec}, \rh(t)\right] -  \tau_c\left[ H_{\scriptscriptstyle \rm sec},\left[H_{\scriptscriptstyle \rm sec},\rh(t) \right] \right],\quad [0<t<\tau_1]
\end{eqnarray}
 Here, $H_{\scriptscriptstyle \rm sec}$ is the interaction representation of the secular part of $\mathcal{H}_x(t) + \hdd$. The form of the secular Hamiltonian is given by, $H_x = \sum\limits_{i=1}^2 \omega_1 I_{x}^i$ and $H_{\scriptscriptstyle \rm dd} = \omega_{d_0} \left(2 I_z^1 I_z^2 - I_x^1 I_x^2 - I_y^1 I_y^2 \right) $. In general, for a two-spin system, $\rh$ can be written by using fifteen observables. The above observables can be divided into two parts, nine symmetric observables, and six anti-symmetric observables. As the initial states are a linear combination of the symmetric observables, our dynamics can be described by only those nine observables. The form of the  symmetric observables is given below,
\begin{eqnarray}\label{obs}
\rh (t) &=& \sum\limits_{\alpha, \beta }  A_{\alpha \beta} (t)I_{\alpha} \otimes I_{\beta},
\end{eqnarray}
where, $\alpha$, $\beta \in \{x,y,z,d\}$, and $I_d = 2\times 2$ identity matrix.  $M_i = A_{id}  + A_{di}$, $M_{ii} = A_{ii}$, $M_{ij} = A_{ij} + A_{ji}$ $\forall i\neq j$ and $i,j \neq d$. Recently Chakrabarti \etal and Saha \etal showed the existence of prethermal plateau in dipolar systems using the above dynamical equation (Eq. (\ref{dynamics-1})) \cite{chakrabarti_2022_creation,saha_cascaded_2023}. In terms of observables, the dynamics can be divided into two subgroups $\{M_x,\, M_{yy}, M_{zz}, M_{yz}\}$ and  $\{M_z,\, M_{y}, M_{xz}, M_{xy}\}$. The dynamics are constrained by several conserved quantities. Such quantities can be written using the linear combination of the observables from the first subgroup. Their analytical forms are given as $3\omega_{d_0}\dot{M}_{zz}+\omega_1  \dot{M}_{x} = 0$,  
$\dot{M}_{yy}+\dot{M}_{zz} =0$ and $\dot{M}_{xx} = 0$ \cite{saha_cascaded_2023}. The solution of the above Eq. (\ref{dynamics-1}), in terms of the observable is denoted as $M_i^{\scriptscriptstyle \rm pre}(\alpha, t)$. Here `pre' denotes the prethermal state and $\alpha$ is the initial value of the observable.  The value of $M_x(t)$ in the prethermal state is written as, 
\begin{eqnarray}
M_x^{\scriptscriptstyle \rm pre}(M_\circ,t_{\scriptscriptstyle \rm pre})  = M_\circ \frac{\omega_1^2}{\omega_1^2 + \frac{9}{16}\omega_{d_0}^2}
\label{mxpre}
\end{eqnarray}
Here, $t_{\scriptscriptstyle \rm pre}$, denotes the time required to reach the prethermal state.
In experiments, drive strength is taken to be stronger than the dipolar interaction ($\omega_1 > \omega_{d_0}$), such consideration leads to the negligible decay of $M_x(t)$ in the transient phase \cite{choi_observation_2017}. 

After the spin-locking sequence, a $\pi+ \delta$ rotation is given along the `$y$' direction for a time period $\tau_2$. No drive-induced dissipation occurs at this time as $\tau_2$ is small. \cite{beatrez_critical_2023}. The dynamical equation for this case is written as,
\begin{eqnarray}\label{dynamics-2}
\frac{d\rh}{dt}&=& -i\left[H_y, \rh(t)\right], \qquad [\tau_1<t<\tau_2].
\end{eqnarray}
Here, $H_y = \omega_2 \sum_{i=1}^2 I_y^i$ and $\omega_2\tau_2 = \pi + \delta$. For the initial condition, $M_x \stoo = M_\circ$, the solution of the above dynamical equation (Eq. \ref{dynamics-2}) is given by, $M_x^{\scriptscriptstyle \rm rot} (t) = M_\circ \cos \omega_2 t$ and $M_z^{\scriptscriptstyle \rm rot} (t) = M_\circ \sin \omega_2 t$. To obtain the $2\tau$ periodic response in the system, the total sequence, $\tau = \tau_1 + \tau_2$ is repeated for $n$ times $[n \in \mathcal{I}]$. For the numerical simulation, it is convenient to analyze the dynamical equation [Eq. (\ref{dynamics-1},\ref{dynamics-2})] in the Liouville space. The dynamical equation in this space is written as $\dot{\hat{\rh}}= \hat{\mathcal{L}}\hat{\rh}$. In this case, $\hat{\mathcal{L}}$ is a $16 \times 16$ matrix, whereas, $\hat{\rh}$ is a $16 \times 1$ column matrix. The solution of $\hat{\rh}(t)$ after the $n$-cycle is written as,
\begin{eqnarray}\label{main-eq}
\hat{\rh}(t) =\left[e^{\hat{\mathcal{L}}_y  \tau_2} e^{\hat{\mathcal{L}}_{\scriptscriptstyle \rm sec}  \tau_1}\right]^n \hat{\rh} \stoo.
\end{eqnarray}
Here, $\hat{\mathcal{L}}_y$ is the Liouvillian corresponding to $H_y$. We also define, $\hat{\mathcal{L}}_1$  is the Liouvillian corresponding to $H_{\scriptscriptstyle \rm dd} + H_x$, therefore, $\hat{\mathcal{L}}_{\scriptscriptstyle \rm sec} = \hat{\mathcal{L}}_1 + \tau_c \hat{\mathcal{L}}_1 \times\hat{\mathcal{L}}_1 $.


\begin{figure*}[htb]
\includegraphics[scale =0.75]{Figure-2.pdf} 
\caption{Plot of $M_x$ versus no. of cycle is shown in  (a), (c),  (e) and their corresponding Fourier transform, $S(M_x)$ versus $\omega$ are shown in  (b), (d), and (f). The value of the fixed parameters are given as, $\omega_1 = 2\pi \times 50$ KHz, $\omega_2 = 2\pi \times 100$ KHz, $\omega_{d_0} = 2\pi \times 2$ KHz, $\tau_c = 10^{-3}$ msec, and $n=200$. The initial state is chosen as, $\vert\psi\rangle \stoo = \prod_{i=1}^2 \left(\vert\uparrow \rangle_i + \vert\downarrow \rangle_i\right)/ \sqrt{2}$. For the upper plot (Fig. 2 (a), (b)), the value of $\tau_1 = 2\pi\times 0.02/ \omega_1$ KHz$^{-1}$, $\tau_2 = \pi/ \omega_2$ KHz$^{-1}$. For the middle plot (Fig. 2 (c), (d)), the value of $\tau_1 = 2\pi\times 0.02/ \omega_1$ KHz$^{-1}$, $\tau_2 = 1.04\pi/ \omega_2$ KHz$^{-1}$. For the lower plot (Fig. 2 (e),(f)), the value of $\tau_1 = 2\pi/ \omega_1$ KHz$^{-1}$, $\tau_2 = 1.04\pi/ \omega_2$ KHz$^{-1}$. The period-doubling occurs for the upper plot (Fig. 2(b)) as the peak of the Fourier spectra appears at $\omega = \pi$. A small increase in $\omega_2 \tau_2$ destroys the $2\tau$ periodic response in (shown in Fig. 2(d)), which can be retrieved for a significant increment of $\tau_1$ (shown in Fig. 2(f)). }
\label{fig-2}
\end{figure*}
 We numerically solve Eq. (\ref{main-eq}) and calculate $M_x(t)$ and its spectrum for the initial condition $M_x\stoo = 1$ and plot them in Fig. \ref{fig-2}. For a fixed $\omega_{d_0},\, \omega_1,\, n$ and $\tau_c$, when $\omega_2 \tau_2 = \pi$ and $\tau_1 = 2\pi\times 0.02/ \omega_1 $ KHz, we observe a $2\tau$ periodic response. For $\omega_2 \tau_2  = \pi$, the `$x$' magnetization flipped from $\hat{x}$ to $-\hat{x}$ in each cycle, therefore the spectrum has a single peak at a position $\omega = \pi$ (upper curve in Fig. \ref{fig-1}). Next, we study the stability of the $2\tau$ response by varying the experimental parameters $\tau_1$ and $\delta$. For the same $\tau_1$, if $\delta = 0.04 \pi$, the $2\tau$ periodicity has vanished and we get two peaks very close to $\omega = \pi$.  Such period-doubling response can be retrieved with a larger $\tau_1$ (100 times than the previous one).  Therefore, our result suggests that lower $\theta$ and higher $\tau_1$ are desirable for observing the DTC phase. Following the works by Choi \etal, we define the crystalline fraction ($f$) as, $f = \vert S(\omega = 0.5)\vert^2 / \sum_\omega \vert S(\omega )\vert^2$ \cite{choi_observation_2017}. We also show a contour plot of crystalline fraction ($f$) as a function of $\omega_2 \tau_2 $ and $\omega_1 \tau_1$  to capture the regime of the DTC phase in the system. The plot of $f$ clearly shows the dependence of $\delta$ and $\tau_1$  on the DTC phase. The stable regime is shown in the yellow color of the contour plot in Fig. 3a,  which also matches with earlier experimental results \cite{choi_observation_2017}. 
 \begin{figure*}[htb]
 %\hspace*{-1mm}
 \raisebox{3cm}{\normalsize{(a)}}\hspace*{-1mm}
\includegraphics[width=0.36\linewidth]{Figure-3.pdf} 
\hspace*{2mm}
\raisebox{3cm}{\normalsize{(b)}}\hspace*{-1mm}
\includegraphics[width=0.36\linewidth]{Figure-4.pdf} 
\raisebox{3cm}{\normalsize{(c)}}\hspace*{-1mm}
\includegraphics[width=0.36\linewidth]{Figure-5.pdf} 
\hspace*{2mm}
\raisebox{3cm}{\normalsize{(d)}}\hspace*{-1mm}
\includegraphics[width=0.36\linewidth]{Figure-6.pdf}
\caption{Fig. (a) shows the contour plot of the crystalline fraction $(f)$ as a function of $\omega_1 \tau_1$ and $\omega_2 \tau_2$. The list of fixed parameters are given here, $\omega_1 = 2\pi \times 20$ KHz, $\omega_2 = 2\pi \times 1000$ KHz, $\omega_{d_0} = 2\pi \times 0.2$ KHz, $\tau_c = 10^{-4}$ msec, and $n=200$. We note that the DTC phase is more robust for lower values of $\omega_2 \tau_2$ and higher values of $\omega_1 \tau_1$. Fig. (b) shows the contour plot of $M_x$ as a function of $\omega_2 \tau_2$ and no. of  cycle. The list of fixed parameters are given here, $\omega_1 = 2\pi \times  20$ KHz, $\omega_2 = 2\pi \times 1000$ KHz, $\omega_{d_0} = 2\pi \times 2$ KHz, $\tau_c = 10^{-3}$ msec, $\tau_1 = 2\pi/ \omega_1$ KHz$^{-1}$ , and $n=50$. The blue region around $\omega_2 \tau_2 = 0$ is known as the prethermal phase. The alternative blue and red stripe around $\omega_2 \tau_2 = \pm \pi$, which shows the existence of the DTC phase.   Fig. (c) shows the contour plot of $\vert S(M_x) \vert^2$ as a function of $\omega_2 \tau_2$ and $\omega$. The list of fixed parameters are given here, $\omega_1 = 2\pi \times  20$ KHz, $\omega_2 = 2\pi \times 1000$ KHz, $\omega_{d_0} = 2\pi \times 0.5$ KHz, $\tau_c = 10^{-3}$ msec, $\tau_1 = 2\pi \times 2 / \omega_1$ KHz$^{-1}$ , and $n=200$. The Fourier peak around $\omega = \{0,\, 2\pi\}$, arises due to the presence of the prethermal phase, similarly, the peaks around $\omega = \pi$ denote the emergence of the DTC phase. Fig. (d) shows the contour plot of $\vert S(M_x) \vert^2$ as a function of $\tau_c$ and $\omega$. The list of fixed parameters are given here, $\omega_1 = 2\pi \times  50$ KHz, $\omega_2 = 2\pi \times 1000$ KHz, $\omega_{d_0} = 2\pi \times 0.1$ KHz, $\tau_1 = 2\pi \times 2 / \omega_1$ KHz$^{-1}$, $\tau_2 = \pi \times 1.01/ \omega_1$ KHz$^{-1}$ , and $n=200$. We find that the DTC phase is more robust in the high $\tau_c$ regime. Lowering $\tau_c$ leads to the breaking of $2\tau$ periodicity in the system.}
\label{fig-3}
\end{figure*}

For $\omega_2\tau_2 = 0$, the solution of the dynamical equation (Eq. (\ref{main-eq})) shows the prethermalization in the system as $M_x (t)$ is constant (Eq. (\ref{mxpre})). Therefore, in Fig. \ref{fig-3}(b), prethermalization occurs in the blue regime (color online) near the $\omega_2\tau_2 = 0$. On the other hand, the system reaches the DTC phase around  $\omega_2\tau_2 = \pm \pi$. The successive red and blue stripes (color online)  occur due to the flipping of $M_x$ at each cycle, $[(M_x)_{2n +1} = -(M_x)_{2n}]$, which shows the notion of a robust $2\tau$ periodicity around $\omega_2\tau_2 = \pm \pi$.

 To confirm the existence of the DTC phase in the system, we also provide a contour plot of the spectrum of $M_x$ in Fig. \ref{fig-3}(c).  The prethermal phase occurs when $\omega_2\tau_2 = 2N \pi$, $\{N \in I\}$. Therefore, the Fourier peak arises at $ \omega = 0, \, 2\pi$ for that particular values of $\omega_2 \tau_2$, which signifies that the applied drive and the response have the same periodicity at the prethermal phase. From Fig. \ref{fig-3}(b), we find that, DTC phase occurs at $\omega_2 \tau_2 = \pm \pi$, similarly the Fourier peaks in Fig. \ref{fig-3}(c), arises at $\omega =  \pi$ for $\omega_2 \tau_2 \approx \pm (2N +1) \pi$. Such period doubling exhibits the notion of the DTC phase. The results corresponding to the Fig. \ref{fig-3}(b), (c) are in agreement with the experimental observation by Beatrez \etal \citep{beatrez_critical_2023}. 

In our dynamical equation (Eq. (\ref{main-eq})), $\tau_c$ plays an important role. To demonstrate the $\tau_c$ dependence of the DTC phase, we provide a simplified analytical calculation. Here, we assume that $\delta/ \pi \ll 1 $, so, $\sin(\pi +\delta) \approx -\delta, \, \cos(\pi +\delta)  \approx 1$. In table \ref{table-1}, the solution of the dynamical equation (Eq. (\ref{main-eq})) is provided for each time instance $\tau_1, \, \tau_2$ up to the period $2\tau$. Here, $\omega_2 \tau_2 = \pi + \delta$.
\begin{table}[htb]
\begin{center}
\caption{Solution of $M_x (t)$ at every time instances for $M_x (t)$.}
\begin{tabular}{ p{1.5cm} |p{4.5cm}    }
 Time$(t)$& $M_x (t)$\\
 \hline
 \hline
 $0$ & $M_{\circ}$\\
 $\tau_1$  &  $M_x^{\scriptscriptstyle \rm pre}(M_\circ,\tau_1)$    \\
 $\tau $  & $-M_x^{\scriptscriptstyle \rm pre}(M_\circ,\tau_1)$     \\
 $\tau + \tau_1$  &  $-M_x^{\scriptscriptstyle \rm pre}(M_x^{\scriptscriptstyle \rm pre}(M_{\circ}, \tau_1), \tau_1)$    \\
 $2\tau$  &  $ M_x^{\scriptscriptstyle \rm pre}(M_x^{\scriptscriptstyle \rm pre}(M_{\circ}, \tau_1), \tau_1) +$    \\
 &   $\delta^2 M_z^{\scriptscriptstyle \rm pre}(M_x^{\scriptscriptstyle \rm pre}(M_{\circ}, \tau_1), \tau_1)$ \\
 \hline
  \end{tabular}
\label{table-1}
\end{center}
\end{table} 
The solution shows that the `$x$' magnetization is flipped in every first cycle of time-period $\tau$. In the next cycle, $M_x$ doesn't flip perfectly, as $\delta \neq 0$. Therefore, a small imperfection in the $\pi$ rotation along the `$y$' direction destroys the DTC order in the system. Although, a sufficiently long $\tau_1$ helps to bring back the DTC order because for large $\tau_1$. The extra term of $M_x (2\tau)$ in the table \ref{table-1} decays to zero as in the limits of $\omega_1> \omega_{d_0}$, the analytical form of the extra term is written as, $M_z^{\scriptscriptstyle \rm pre}(\alpha, \tau_1) \approx \alpha\cos(\omega_1 \tau_1) e^{-\omega_1^2 \tau_1 \tau_c}$. Hence the solution becomes, $M_x (2\tau) = M_x^{\scriptscriptstyle \rm pre}(M_x^{\scriptscriptstyle \rm pre}(1, \tau_1), \tau_1) $.  In the above case, $2\tau$ periodicity can be also brought back by changing $\tau_c$ as well.  We find that the role of $\tau_1$ and $\tau_c$ are complementary to each other. It is well known that, $\tau_c$ is inversely proportional to the temperature \cite{chakrabarti2018a}. Hence, by lowering the temperature the effect of the imperfect rotation can be diminished and the sub-harmonic response can be restored. We also show a contour plot of $\vert S(M_x)\vert^2$ as a function of $\tau_c$ and $\omega$, which is equivalent to the analytical expression, which is provided in table \ref{table-1}. Therefore, the $2\tau$ periodic response vanishes for lowering $\tau_c$. 

We note that, the $2\tau$ periodicity can occur for $\omega_1\tau_1 = (2N +1)\pi/2$. Although, such a response is not a DTC order and it is previously reported as the appearance of accidental dynamical decoupling, which looks similar to the $2\tau$ periodicity of the DTC phase \cite{choi_observation_2017}. Therefore, we do not consider such points in our numerical simulations.

In the $\tau_1$ time duration, the dynamics is governed by both first-order unitary and second-order dissipative processes. Such dissipation is essential for the persistence of DTC order in the system. Due to the second-order terms, the extra effects due to the imperfect rotation in $\tau_2$ time duration are reduced. In addition to that, $M_x$ is stabilized by the dissipation,  as the dissipation doesn't destroy the conserved quantity  $(\omega_1\dot{M}_x + 3 \omega_{d_0}\dot{M}_{zz}  = 0)$,  which helps to retain the DTC phase. Such novel effects cannot be captured by using the formal form of the QME. Thus, only FRQME can describe the origin of DTC in a dissipative dipolar system. In our case, the effects of $\hsl$ and non-secular terms of $\hdd$ are neglected, as they provide a very slow decay rate. In Fig. \ref{fig-3}(a),  we show that for higher $\omega_1 \tau_1$, the DTC phase is stabilized. Although, in real experiments, the DTC regime becomes narrower for longer values of $\omega_1 \tau_1$ \cite{choi_observation_2017,ho_critical_2017}. Such narrowing occurs due to the presence of other sources of slow decays which we have neglected in this case. We note that the existence of the dissipation stabilized DTC is recently reported by other groups in a different context \cite{lazarides_time_2020}.



As a summary, we present an alternative approach to describe the emergence of the DTC phase in a dissipative dipolar system under the two-drive protocol, which has been used in recent experiments \cite{choi_observation_2017,beatrez_critical_2023}. We consider a two-spin dipolar-coupled dissipative system. Using FRQME, we obtain the existence of the DTC phase and numerically plot the phase diagram in terms of the experimental parameters. Our results are well matched with the recent experiments\cite{choi_observation_2017,beatrez_critical_2023}. Another important finding is the dependence of the DTC phase on the environmental parameter $(\tau_c)$, which suggests that the DTC phase is more stable at low temperatures.

Our analysis is based on a simplified dynamical model consisting of two-spin. In general, a prescription considering a many-spin dipolar coupled network is required for exact matching with the experimental results. We also neglect the decay of the magnetization modes in the period $t_{acq}$, which was considered in the experiments by Beatrez \etal \cite{beatrez_critical_2023}. Finally, our dynamical equation is not effective in very low temperatures (high $\tau_c$), as there is no clear separation between the system and environmental time-scale in that temperature regime. In such cases, a non-Markovian QME may be useful but further investigation is needed.   










%\begin{acknowledgments}

The authors thank Arpan Chatterjee for their insightful comments and helpful
suggestions. SS gratefully acknowledges University Grants Commission (UGC) of Govt. of India for a research fellowship (Student ID: MAY2018- 528071).

%\end{acknowledgments}
 
\bibliographystyle{apsrev4-2}
\bibliography{references}
\end{document}








